$\mathtt{CUT} \ge \log_2 m$ is that $a^{\phi(m)} \not\equiv 1 \pmod m$ can happen. But for $d \ge \log_2 m$, $a^{d+\phi(m)} \equiv a^{d} \pmod m$.
\begin{minted}{cpp}
const ll INF = 1e18; // INF > mod^2; might need i128
ll fpow(ll mul, ll bit, ll mod, ll lim = INF){
	ll res = 1; while (bit){
		if (bit&1){ res = min(res*mul, lim) % mod; }
		mul = mul*mul % mod; bit >>= 1;
	} return res;
}

const int CUT = 100, LEN = 4; // CUT >= log2(mod); A = [2 for 1..LEN] > CUT
int n; int arr[1000020];
int solve(int idx){
	int val = 1; for (int i = n; i >= idx; i--){
		val = fpow(arr[i], val, (CUT+1)*(CUT+2), CUT+1);
	} return val;
}
int solve(int idx, int mod){
	if (mod == 1){ return 0; }
	if (idx == n){ return arr[idx] % mod; }
	if (n-idx+1 <= LEN){
		if (solve(idx+1) <= CUT){ return fpow(arr[idx], solve(idx+1), mod); }
	}
	int m = phi[mod];
	return fpow(arr[idx], solve(idx+1, m) + CUT*m, mod);
}
\end{minted}